%% ma: L12-B05-0025
%% lop: 12
%% chuong: 1
%% bai: 5
%% dang: 12.5.17
%% loai: TLN
%% muc: VD
%% dap_an: "5,66"
%% nguon: "(NBV)-ĐỀ SỐ 14-MỖI NGÀY 1 ĐỀ THI 2025.pdf (D14, MỖI NGÀY 1 ĐỀ - Đề số 14), Phần III Câu 5"
%% kien_thuc: "Bài 5 (ứng dụng đạo hàm vào thực tiễn: tốc độ biến thiên), Bài 32 lớp 11 (đạo hàm hàm hợp)"
%% trang_thai: chinh-thuc
%% ghi_chu: "Đã duyệt 15/10/2026. Đáp án gốc đúng (5,66). Ảnh bức tường và chiếc thang ở đề gốc là ảnh chụp trang trí, không có số liệu toán nên bỏ hình."
\begin{ex}
Một chiếc thang dài $9$ mét tựa vào một bức tường thẳng đứng trên một mặt đất bằng phẳng. Khi đầu dưới của thang di chuyển (trên mặt đất) ra xa bức tường với vận tốc không đổi là $2$ (m/s) thì đầu trên cùng của thang sẽ trượt xuống dọc theo bức tường. Khi điểm đầu thang cách mặt đất $3$ mét thì tốc độ di chuyển của nó bằng bao nhiêu? (Đơn vị: m/s, kết quả làm tròn đến hàng phần trăm.)
\shortans{5,66}
\loigiai{Gọi $x(t)$ là khoảng cách từ chân thang đến tường, $y(t)$ là độ cao của đầu trên thang: $x^2+y^2=81$, $y=\sqrt{81-x^2}$, $x'(t)=2$.
$y'(t)=\dfrac{-x\cdot x'(t)}{\sqrt{81-x^2}}=-\dfrac{2x}{y}$.
Khi $y=3$ thì $x=\sqrt{81-9}=6\sqrt2$, nên $y'=-\dfrac{2\cdot6\sqrt2}{3}=-4\sqrt2$.
Tốc độ của đầu thang là $4\sqrt2\approx5{,}66$ (m/s).}
\end{ex}
